\newcommand{\N}{\mathbb{N}}
\newcommand{\R}{\mathbb{R}}
\newcommand{\Q}{\mathbb{Q}}
\newcommand{\Z}{\mathbb{Z}}
\newcommand{\C}{\mathbb{C}}

\DeclareMathOperator{\Jac}{Jac}
\DeclareMathOperator{\Ker}{Ker}
\DeclareMathOperator{\mesh}{mesh}
\DeclareMathOperator{\Var}{Var}
\DeclareMathOperator{\hm}{hm}
\DeclareMathOperator{\Tr}{Tr}

\let\temp\phi
\let\phi\varphi
\let\varphi\temp

\newcommand{\pa}[1]{\left(#1\right)}
\newcommand{\bra}[1]{\left[#1\right]}
\newcommand{\cbra}[1]{\left\{#1\right\}}

\newcommand{\mat}[1]{\begin{matrix}#1\end{matrix}}
\newcommand{\pmat}[1]{\pa{\mat{#1}}}
\newcommand{\bmat}[1]{\bra{\mat{#1}}}

\newcommand{\pfrac}[2]{\pa{\frac{#1}{#2}}}
\newcommand{\bfrac}[2]{\bra{\frac{#1}{#2}}}
\newcommand{\psfrac}[2]{\pa{\sfrac{#1}{#2}}}
\newcommand{\bsfrac}[2]{\bra{\sfrac{#1}{#2}}}

Broadly speaking, we will consider problems of the form \[
    \min_{x \in \mathcal X} f(x)
\] for (potentially) randomized \(f\) and \(\mathcal X\) sufficiently
large such that brute force is not feasible.

\subsubsection{Examples}\label{examples}

\begin{enumerate}
\def\labelenumi{\arabic{enumi}.}
\tightlist
\item
  Consider the problem of finding the top eigenvalue of a random matrix
  \(G\) in the Gaussian orthogonal ensemble: \[
   \mathbb E \max_{x \in \mathbb{R}^N, |x|_2 = 1} \left\langle x, Gx \right\rangle
  \] with \(G_{ij} = G_{ji} \sim N(0, 1/N)\), and \(G_{ii} = 0\).
  Wigner's semicircle law shows that this is \(\approx 2\) for large
  \(N\). Without the maximum, this problem is simple: \[
   \mathbb E[\left\langle x, Gx \right\rangle] = \left\langle x, \mathbb E[G]x \right\rangle = 0
  \] and \[
  \begin{align*}
   \mathop{\mathrm{Var}}(\left\langle x, Gx \right\rangle) &= \mathbb E \left[ \left( 2\sum_{i < j}x_i G_{ij}x_j \right)^2 \right] \\
   &= 4\mathbb E \left[ \sum_{i < j} x_i^2 G_{ij}^2 x_j^2 \right] \\ 
   &= \frac{4}{N} \sum_{i < j} x_i^2 x_j^2 \\
   &= \frac{2}{N} \left( \left(\sum_{i = 1}^n x_i^2 \right)^2 - \sum_{i=1}^n x_i^4 \right) \\
   &= \frac{2}{N} \left( 1 - \|x\|_4^4 \right),
  \end{align*}
  \] so \(\left\langle x, Gx \right\rangle \sim N(0, 2/N)\).
\item
  Alternatively, consider \[
   \mathbb E \max_{x \in \{\pm 1\}^N} \left\langle x, Gx \right\rangle.
  \] While harder, we can demonstrate that it is bounded above by \(2N\)
  by reducing to the previous case: \[
   \frac{1}{N}\mathbb E \max_{x \in \{\pm 1\}^N} \left\langle x, Gx \right\rangle = 
   \mathbb E \max_{x \in \{\pm 1\}^N} \left\langle \frac{x}{\sqrt{N}}, G \frac{x}{\sqrt{N}} \right\rangle \leq 
   \mathbb E \max_{x \in \mathbb{R}^N, |x|_2 = 1} \left\langle x, Gx \right\rangle \approx 2.
  \]
\item
  (A Random Convex Program) Take a (non-symmetric) random matrix
  \(X \in \mathbb R^{n \times p}\) with \(X_{ij} \sim N(0, 1)\) and \[
   Y = X \beta + \xi
  \] for \(\xi \sim N(0, 1)\). Then we might try to ask, e.g., \[
   \min_{\hat \beta \in \mathbb{R}^p} \|Y - X\hat \beta\|_2^2 + \lambda \|\hat \beta\|_1
  \] in hopes of recovering a sufficiently sparse \(\beta\) vector.
  Being able to compute properties of the achieved minimum can translate
  into computing properties of \(\hat \beta\) (relative to the ``true''
  underlying \(\beta\)).
\item
  (A Nonconvex/Combinatorial Problem) Let
  \(A \in \mathbb R^{N \times N}\) with
  \(A_{ij} \sim \mathrm{Ber}(1/2)\), \(A_{ii} = 1\). Then we might ask
  about \[
   \mathbb E \max \{|S| \mid S \subset [N], A_{SS} = \mathbf 1 \}
  \] which we call the largest clique problem (think of \(A\) as a
  random adjacency matrix for a graph). Further, we might care about
  limitations on ``compute'': what can we do in polynomial time?
\item
  (An Application) One sends radio signals as waves; but there is
  usually some corruption between the sent and received messages. So how
  does one send messages over a noisy channel? The core math problem is
  as follows: given a message \(M \in \{\pm 1\}^{n}\), is there a way to
  encode it into \(\{\pm 1\}^{N}\) such that decoding is robust to
  corruption? For modern solutions, look for low density parity checks:
  write \[
   Y = (M, Z) + \mathrm{ noise}
  \] where \(Z\) is some additional information, and take \(Z_{i}\) to
  be the parity of the sum of a bunch of (well chosen) entries of \(M\).
  Then decoding is a problem that looks like \[
   \hat M = \argmin_{M \in \{\pm 1\}^n} \|Y - \mathrm{enc}(M)\|_0.
  \]
\end{enumerate}

\begin{center}\rule{0.5\linewidth}{0.5pt}\end{center}

\subsection{Gaussian Processes}\label{gaussian-processes}

\emph{Def}: Take some indexing set \(T\); we say that a collection of
random variables\(\{g_t\}_{t \in T}\) is a Gaussian process if all the
finite-dimensional distributions are jointly Gaussian: that is, for all
\(|S| < \infty\), \(g_s = \{g_s\}_{s \in T}\) satisfies
\(g_s \sim N(\mu_S, \Sigma_s)\).
